\section{Conclusão}

O monitoramento ambiental apresenta requisitos compatíveis com o emprego de \acp{RSSF}, como observação simultânea em múltiplos pontos, amostragem periódica e operação prolongada em áreas extensas, remotas ou de difícil acesso. Uma implementação completa reúne sensores selecionados conforme as variáveis ambientais, nós de baixo consumo, processamento e armazenamento local, comunicação sem fio, gateway com conexão de backhaul e uma plataforma capaz de preservar medições, marcas temporais, indicadores de qualidade e informações de calibração. O armazenamento local mantém a continuidade da aquisição durante interrupções temporárias da comunicação.

A partir desses requisitos, a tecnologia de comunicação deve ser selecionada conforme a distribuição dos nós, o volume e a frequência das mensagens, a autonomia requerida, as condições de propagação, a infraestrutura disponível e o custo. LoRaWAN apresenta características compatíveis com nós dispersos e tráfego reduzido, enquanto IEEE 802.15.4 e 6LoWPAN atendem redes locais densas ou com encaminhamento multi-hop. Em áreas com cobertura celular, NB-IoT e LTE-M utilizam a infraestrutura da operadora. Wi-Fi, por sua vez, atende estações alimentadas continuamente ou aplicações com maior volume de dados. A configuração selecionada requer ensaios de cobertura, perda de pacotes, margem do enlace e consumo energético no local da implantação.

A qualidade científica da RSSF depende da instalação correta dos sensores, da calibração, da consistência temporal, da representatividade espacial, da continuidade das séries e da viabilidade da manutenção. A implantação piloto permite avaliar esses requisitos durante um ciclo ambiental representativo antes da expansão da rede. O resultado esperado é uma série ambiental confiável, rastreável e adequada à finalidade definida para o monitoramento.